\bookchapter{제44장: 0에서 1까지 실전 — 완성 프로젝트 전 과정}

\chapterintro{이론은 끝났으니 이제 손을 움직일 시간이에요. 이번 장은 단순히 “따라 치기” 강좌가 아니라 “어떻게 생각하고, 어떻게 묻고, 어떻게 검증할지”에 대한 실전 방법론이에요. 난이도가 다른 세 프로젝트로 요구사항부터 완성품까지의 사고 과정 전체를 보여 드릴게요. 어떤 도구를 쓰든 이 흐름은 그대로 적용돼요.}

\booksection{프로젝트 1: 미니멀 개인 도구 사이트(입문)}

\booksubsection{요구사항 생각하기}

목표: 실용적인 소도구 3\textasciitilde{}4개의 입구를 보여주는 개인 도구 사이트 웹페이지를 만들어요. 순수 프런트엔드 프로젝트라 백엔드와 데이터베이스가 필요 없어요. 중요한 건 도구 자체가 얼마나 복잡한가가 아니라 “요구사항에서 완성품까지”의 전체 흐름을 경험하는 거예요.

요구사항 5요소에 따르면: Who — 개인 사용. What — 도구 입구를 보여주는 웹페이지. Why — 자주 쓰는 도구를 한 페이지에 모아 접근하기 편하게. How — 단일 파일 HTML+CSS+JS, 프레임워크 미사용. Constraint — 30분 안에 완성, 배포 불필요.

\booksubsection{지시 내리기}

Think Before Coding 원칙에 따라, 지시는 AI가 목표를 이해할 만큼 충분한 정보를 담되 모든 구현 세부사항을 지정할 필요는 없어요:

\begin{lstlisting}
개인 도구 모음 웹페이지를 만들어 주세요. 요구사항:
1. 단일 파일 HTML(인라인 CSS·JS), 프레임워크 없이
2. 다크 테마, 현대적인 미니멀 스타일
3. 4개의 도구 입구 카드 표시:
   - 텍스트 요약 도구
   - 글자 수 세기 도구
   - 색상 선택기
   - QR 코드 생성기
4. 각 카드에는 아이콘, 제목, 짧은 설명이 있어요
5. 카드에 호버 애니메이션 효과
6. 반응형 레이아웃, 휴대폰에서도 사용 가능
7. 상단에 제목 표시줄 하나
완성된 HTML 파일을 바로 제시해 주세요.
\end{lstlisting}

\booksubsection{검수와 반복}

AI가 코드를 내놓으면 반복 검증의 3계층 점검으로 검수해요. 첫째 층 기능 검증: 파일이 브라우저에서 열리나요? 카드 4개가 모두 보이나요? 호버 애니메이션은 정상인가요? 둘째 층 경계 검증: 휴대폰 화면 크기에서도 레이아웃이 정상인가요? 창이 아주 좁을 때 카드가 줄바꿈되나요? 셋째 층 아키텍처 검증: 코드가 읽기 쉬운가요? CSS가 style 태그 하나에 인라인되어 있나요? 눈에 띄는 중복은 없나요?

어떤 도구 카드의 기능까지 실제로 구현하고 싶다면(입구만이 아니라) 하나씩 추가하면 돼요. — 도구 기능을 하나 추가하고 검증한 뒤 다음을 추가하는 거예요. 이것이 “작은 단계 빠르게” 전략의 실전 활용이에요.

\booksection{프로젝트 2: 풀스택 가계부 시스템(중급)}

\booksubsection{요구사항 생각하기와 PRD}

목표: 프런트엔드, 백엔드, 데이터베이스를 갖춘 풀스택 가계부 앱. 이 프로젝트는 첫 번째보다 훨씬 복잡해서 제40장의 요구사항 엔지니어링 방법이 필요해요.

먼저 PRD를 써요. 제40장 3절의 템플릿에 따라 프로젝트 개요, 기능 목록, 기술 방안, 데이터 모델, 개발 계획을 작성하세요. 기술 방안 선택은 Simplicity First 원칙을 고려해야 해요. — 프런트엔드는 React를 쓰되 Redux 같은 복잡한 상태 관리는 도입하지 않고, 백엔드는 Node.js를 쓰되 마이크로서비스는 도입하지 않고, 데이터베이스는 (개발 단계에서는) PostgreSQL 대신 SQLite를 쓰는 식이에요.

개발 계획은 5단계로 쪼개요: Step 1 백엔드 API 뼈대(Express 서버 + SQLite + CRUD 인터페이스). Step 2 프런트엔드 기본 페이지(기록 목록 + 추가 폼). Step 3 프런트·백엔드 연동(API 호출 + 데이터 표시). Step 4 분류와 필터 기능. Step 5 월별 요약과 간단한 차트. 매 단계가 끝날 때마다 “쓸 수 있는 것”이 완성돼요.

\booksubsection{프로젝트 기억 파일 설정}

코딩을 시작하기 전에 먼저 프로젝트 기억 파일을 만드세요. 이것이 컨텍스트 엔지니어링의 핵심 실천이에요. — 매 대화가 시작될 때 AI가 프로젝트 배경을 알게 만드는 거예요.

\begin{lstlisting}
# 프로젝트: 풀스택 가계부 앱
## 개요
수입·지출 기록, 분류, 월간 집계를 지원하는 개인 가계부 앱이에요.
## 기술 스택
- 프런트엔드: React + TypeScript + Vite
- 백엔드: Node.js + Express + TypeScript
- 데이터베이스: SQLite(better-sqlite3 사용)
- 프런트·백 공유 타입 정의
## API 규칙
- 기본 경로: /api
- 응답 형식:{ success: boolean, data: any, error: string }
- RESTful 스타일
## 데이터 모델
Record { id, amount, type, category, note, date }
## 코드 규칙
- TypeScript 스트릭트 모드
- 들여쓰기 2칸
- 함수에 JSDoc 주석
- 주석과 커밋 메시지는 한국어로
\end{lstlisting}

\booksubsection{단계별 구현}

개발 계획에 따라 매 단계를 독립된 대화 하나로 완성해요. 다음은 첫 단계의 지시 예시예요:

\begin{lstlisting}
백엔드 API 뼈대를 만들고 싶어요.
요구사항:
1. Express 서버, 포트 3001
2. SQLite 데이터베이스, 테이블 이름 records
3. records 테이블 필드:
   id (TEXT PRIMARY KEY)
   amount (REAL)
   type (TEXT: 'income'|'expense')
   category (TEXT)
   note (TEXT)
   date (TEXT: YYYY-MM-DD)
4. CRUD 엔드포인트:
   GET /api/records - 모든 기록 조회
   POST /api/records - 기록 추가
   PUT /api/records/:id - 기록 수정
   DELETE /api/records/:id - 기록 삭제
5. 에러 처리 미들웨어
완성된 백엔드 코드 구조를 제시해 주세요.
\end{lstlisting}

이 지시의 특징에 주목하세요: 현재 단계(백엔드 뼈대)만 설명하고 프런트엔드 요구는 섞지 않아요. 필드와 인터페이스를 명확히 나열해(AI에게 충분한 제약을 주면서) 구체적 구현은 AI가 결정하게 두죠. 그리고 Goal-Driven Execution 원칙, 즉 “어떻게 하는지”가 아니라 “무엇을 할지”를 말하는 원칙을 따라요.

AI가 코드를 내놓으면 바로 검증해요: 서버를 실행하고 curl이나 Postman으로 각 인터페이스가 정상 작동하는지 테스트하세요. 문제가 없음을 확인한 뒤 새 대화를 열어 Step 2(프런트엔드 기본 페이지)로 넘어가요. 매 단계마다 “지시-검증-새 대화” 순환을 따르는 거예요.

\booksection{프로젝트 3: AI 강화 앱(고급)}

\booksubsection{프로젝트 구상}

목표: 대형 모델 API를 연동한 앱 — 예컨대 “AI 학습 도우미”로, 사용자가 학습 내용을 입력하면 AI가 요약, 관련 지식 포인트, 학습 자료 추천을 생성해요. 이 프로젝트의 특별한 점은: 앱 자체가 AI를 호출하기 때문에 API 호출, 프롬프트 엔지니어링, 오류 처리 같은 새로운 문제를 다뤄야 한다는 거예요.

이 프로젝트가 가르쳐 주는 가장 중요한 교훈은 이거예요: “AI를 호출하는 코드를 AI에게 시켜 쓸 때” 프롬프트 엔지니어링에 각별히 주의해야 해요. 앱이 대형 모델 API를 호출할 때 쓰는 프롬프트 역시 이 가이드에서 논의한 원칙 — 명확함, 간결함, 목표 지향 — 을 따라야 해요.

\booksubsection{핵심 기술 포인트}

이 프로젝트에는 앞선 두 프로젝트에서 만나지 못한 기술 포인트가 몇 가지 등장해요:

\begin{itemize}
  \item API 키 관리: 키는 절대 코드에 하드코딩하지 말고 환경 변수로 관리
  \item 스트리밍 응답 처리: 대형 모델의 답변은 보통 스트리밍 방식이라 한 글자씩 표시해야 함
  \item 토큰 사용량 통제: 대형 모델은 토큰 기준 과금이라 앱 차원에서 사용량을 통제해야 함
  \item 프롬프트 템플릿 관리: 기능마다 다른 프롬프트 템플릿을 쓰므로 체계적으로 관리해야 함
\end{itemize}

이 기술 포인트들은 모두 AI에게 시켜 구현할 수 있지만, 각각의 “왜”를 이해해야 해요. — 왜 키는 하드코딩하면 안 될까? 왜 스트리밍 응답이 필요할까? 왜 토큰 사용량을 통제해야 할까? “왜”를 이해해야 AI의 방안이 합리적인지 판단할 수 있어요.

\booksection{범용 실전 흐름 체크리스트}

프로젝트 크기와 관계없이 쓸 수 있는 범용 AI 코딩 실전 흐름 체크리스트예요. 새 프로젝트를 시작할 때마다 이 체크리스트를 훑으며 핵심 단계를 빠뜨리지 않았는지 확인하세요.

\begin{center}
\begin{tabularx}{\linewidth}{|X|X|X|}
\hline
\textbf{단계} & \textbf{할 일} & \textbf{관련 원칙} \tabularnewline
\hline
1. 요구사항 생각하기 & 요구사항 5요소를 훑으며 무엇을 만들지 분명히 하기 & Think Before Coding \tabularnewline
2. 요구사항 분해 & 큰 요구를 독립적으로 검증 가능한 작은 단계로 쪼개기 & Simplicity First \tabularnewline
3. PRD 작성 & 구조화된 문서로 정리해 참고 자료로 삼기 & Goal-Driven \tabularnewline
4. 기억 파일 설정 & CLAUDE.md 등 동류 파일 만들기 & Harness Engineering \tabularnewline
5. 단계별 구현 & 단계마다 대화 하나, 지시-검증-다음 단계 & Surgical Changes \tabularnewline
6. 테스트 검증 & 단계가 끝날 때마다 즉시 기능 정확성 테스트 & Iterate \& Verify \tabularnewline
7. 코드 리뷰 & 아키텍처 타당성, 보안성, 유지보수성 점검 & Iterate \& Verify \tabularnewline
8. 반복 개선 & 테스트와 리뷰 결과를 반영해 반복해서 개선 & Iterate \& Verify \tabularnewline
\hline
\end{tabularx}
\end{center}

\begin{center}\small\color{gray}AI 코딩 범용 실전 흐름 체크리스트\end{center}

\begin{tcbpractice}{직접 해보기: 이번 주에 “프로젝트 1” 완성하기}
제44장 1절의 “개인 도구 사이트”를 여러분의 첫 AI 코딩 숙제로 삼아 보세요: 5분을 들여 5요소를 훑고, 지시를 AI에 보낸 뒤(어떤 도구든 좋아요) 3계층 검증으로 검수해요. 30분 안에 완성하세요 — 여러분의 첫 AI 코딩 작품이 탄생해요. 다 만들면 꼭 잘 저장해 두세요. “독자”에서 “실천자”로 넘어간 기념품이거든요.

\end{tcbpractice}

\begin{tcbsummary}{이번 장 정리}
\begin{itemize}
  \item 프로젝트 1(입문): 순수 프런트엔드 도구 사이트 — 전체 흐름 체험, 30분이면 완성 가능
  \item 프로젝트 2(중급): 풀스택 가계부 시스템 — PRD + 기억 파일 + 5단계 단계별 구현, 매 단계 “쓸 수 있는” 상태
  \item 프로젝트 3(고급): AI 학습 도우미 — AI에게 AI 호출 코드 작성시키기, 프롬프트 엔지니어링 주의
  \item 범용 흐름 8단계: 요구사항 생각하기 분해 PRD 기억 파일 단계별 구현 테스트 리뷰 반복
  \item 매 단계마다 “지시-검증-새 대화” 순환 적용
\end{itemize}

\end{tcbsummary}
